% SPDX-License-Identifier: CC-BY-SA-4.0
% Part: Syntactic analysis
% Section: Algebraic grammars
% Exercise: Grammar of the calculator

\begingroup

\begin{exercice}[Grammaire de la calculatrice]\label{exo:parsing/grammars/calculator}

  Soit la grammaire :
  $\begin{array}[t]{ll}
    G \eqdef  \langle &  \{a, b, c, +, \times, (, )\}, \{S, \mathit{somme}, \mathit{produit}, \mathit{facteur}, \mathit{terme}\}, S,\\
    & \left\{ \begin{array}{lll}
      S                &\rightarrow& \mathit{somme} \\
      \mathit{somme}   &\rightarrow& \mathit{somme} + \mathit{produit} \mid \mathit{produit} \\
      \mathit{produit} &\rightarrow& \mathit{produit} \times \mathit{facteur} \mid  \mathit{facteur} \\
      \mathit{facteur} &\rightarrow& ( \mathit{somme} ) \mid \mathit{terme} \\
      \mathit{terme}   &\rightarrow& a \mid  b \mid  c
    \end{array} \right\}\;\rangle
  \end{array}$

  \begin{question}
  \item Décrire la structure arborescente (arbre de la syntaxe abstraite) de l'expression $a \times b \times c + a \times (b + c)$.
  \item Donner une génération de la chaîne $a \times b \times c + a \times (b + c)$ par la grammaire $G$.
  \item Justifier la non-appartenance de la chaîne $a + ()$ au langage engendré par $G$.
  \item Modifier la grammaire pour qu'elle prenne en compte les opérations $-$ et $/$,
    dont les propriétés sont respectivement égales à celles de l'addition et du produit.
    Pour cette nouvelle grammaire, donner une génération de la chaîne $a/(d + e - f)$
  \end{question}

\end{exercice}

\begin{correction}

  \begin{question}

  \item La figure suivante montre l'AST associé à la chaîne $a \times b \times c + a \times (b + c)$.
    $$\begin{tikzpicture}[tree]
    \tree{$+$}{
      \tree{$\times$}{
        \tree{$\times$}{
          \tree{$a$}{}
          \tree{$b$}{}
        }
        \tree{$c$}{}
      }
      \tree{$\times$}{
        \tree{$a$}{}
        \tree{$+$}{
          \tree{$b$}{}
          \tree{$c$}{}
        }
      }
    }
  \end{tikzpicture}$$
    
  \item La figure suivante montre l'arbre de dérivation associé à la chaîne $a \times b \times c + a \times (b + c)$.
    $$\begin{tikzpicture}[tree]
    \tree{$S$}{\tree{$\mathit{somme}$}{
        \tree{$\mathit{somme}$}{\tree{$\mathit{produit}$}{
            \tree{$\mathit{produit}$}{
              \tree{$\mathit{produit}$}{\tree{$\mathit{facteur}$}{\tree{$\mathit{terme}$}{\tree{$a$}{}}}
              }
              \tree{$\times$}{}
              \tree{$\mathit{facteur}$}{\tree{$\mathit{terme}$}{\tree{$b$}{}}}
            }
            \tree[xshift=-.5]{$\times$}{}
            \tree{$\mathit{facteur}$}{\tree{$\mathit{terme}$}{\tree{$c$}{}}}
        }}
        \tree{$+$}{}
        \tree{$\mathit{produit}$}{
          \tree{$\mathit{produit}$}{\tree{$\mathit{facteur}$}{\tree{$\mathit{terme}$}{\tree{$a$}{}}}}
          \tree[xshift=.5]{$\times$}{}
          \tree{$\mathit{facteur}$}{
            \tree[xshift=-.5]{$($}{}
            \tree{$\mathit{somme}$}{
              \tree{$\mathit{somme}$}{\tree{$\mathit{produit}$}{\tree{$\mathit{facteur}$}{\tree{$\mathit{terme}$}{\tree[xshift=-1]{$b$}{}}}}}
              \tree[xshift=-1]{$+$}{}
              \tree{$\mathit{produit}$}{\tree{$\mathit{facteur}$}{\tree{$\mathit{terme}$}{\tree[xshift=-1]{$c$}{}}}}
            }
            \tree[xshift=-1.5]{$)$}{}
          }
        }
    }}
  \end{tikzpicture}$$

  \item La règle qui crée des parenthèses crée aussi le 
    non terminal $somme$ qui ne peut pas résulter en $\varepsilon$.

  \item
  $\begin{array}[t]{ll}
    G \eqdef  \langle &  \{a, b, c, +, \times, (, )\}, \{S, \mathit{somme}, \mathit{produit}, \mathit{facteur}, \mathit{terme}\}, S,\\
    & \left\{ \begin{array}{lll}
      S                &\rightarrow& \mathit{somme} \\
      \mathit{somme}   &\rightarrow& \mathit{somme} + \mathit{produit} \mid \mathit{somme} - \mathit{produit} \mid \mathit{produit} \\
      \mathit{produit} &\rightarrow& \mathit{produit} \times \mathit{facteur} \mid \mathit{produit} / \mathit{facteur} \mid \mathit{facteur} \\
      \mathit{facteur} &\rightarrow& ( \mathit{somme} ) \mid \mathit{terme} \\
      \mathit{terme}   &\rightarrow& a \mid  b \mid  c
    \end{array} \right\}\;\rangle
  \end{array}$

    $$\begin{tikzpicture}[tree]
    \tree{$S$}{\tree{$\mathit{somme}$}{\tree{$\mathit{produit}$}{
          \tree{$\mathit{produit}$}{
            \tree{$\mathit{facteur}$}{\tree{$\mathit{terme}$}{\tree[xshift=1]{$\mathit{a}$}{}}}
          }
          \tree[xshift=2.25]{$/$}{}
          \tree{$\mathit{facteur}$}{
            \tree[xshift=2]{$($}{}
            \tree{$\mathit{somme}$}{
              \tree{$\mathit{somme}$}{
                \tree{$\mathit{somme}$}{\tree{$\mathit{produit}$}{\tree{$\mathit{facteur}$}{\tree{$\mathit{terme}$}{\tree{$e$}{}}}}}
                \tree{$+$}{}
                \tree{$\mathit{produit}$}{\tree{$\mathit{facteur}$}{\tree{$\mathit{terme}$}{\tree{$e$}{}}}}
              }
            \tree[xshift=-.5]{$-$}{}
            \tree{$\mathit{produit}$}{\tree{$\mathit{facteur}$}{\tree{$\mathit{terme}$}{\tree{$f$}{}}}}
            }
            \tree[xshift=-1]{$)$}{}
          }
    }}}
  \end{tikzpicture}$$

    
  \end{question}

\end{correction}

\endgroup
\endinput
